%% text-methods.tex -- shared body text: how the margin is drawn
\setcounter{section}{3}\setcounter{subsection}{1}
\subsection{The garden in the margin}

Everything in the left margin of this page is drawn while the page is being shipped out,
from a few hundred lines of Lua run by Lua\TeX. There are no image files, no imported
paths and no external programs; every petal is a parametric curve evaluated for this page
and written straight into it as PDF path operators. This section records how it is done,
so that the margin documents itself.

\begin{definition}
A \emph{primitive} is a closed outline $\gamma\colon[0,1]\to\mathbb{R}^2$ authored by hand
together with a small set of free parameters $(p_1,\dots,p_k)$ and a recipe for painting
it. An \emph{instance} of the primitive samples the parameters from fixed intervals and
paints $\gamma$ in a local frame chosen by the caller. The composition on this page is
\emph{authored}; only the parameters are \emph{sampled}. This is what ``procedural'' means
here: no flower is grown from a grammar, but no two petals share a shape.
\end{definition}

\begin{construction}[Palette]
There is one document color, \texttt{ink}, and from it a ramp $t\mapsto\operatorname{tone}(t)$
from the darkest navy at $t=0$ to a pale mist at $t=1$. The ramp is interpolated in
OKLCH between six stops whose lightness runs from $0.25$ to $0.92$ and whose chroma peaks
at $0.175$ for the cobalt at $t=0.38$; the hue is that of \texttt{ink}, drifting $6^\circ$
toward violet in the darks and $8^\circ$ toward cyan in the lights, as ultramarine does when
it is thinned with water, and the chroma scales with that of \texttt{ink}. Mixing the
document color with white, as \texttt{ink!22} would, desaturates toward gray instead;
the perceptual ramp keeps the light end blue. A second color, \texttt{accent}, is a copper
used only for hairlines.
\end{construction}

\begin{construction}[Petal]
A petal of length $1$ pointing along $+y$ has a spine $x=b\,v^2$, $y=v$ ($b$ bends the
tip) and the obovate half-width
\[
  w(v)=\begin{cases}
    W\bigl(0.16+0.84\,\sin^{1.3}\frac{\pi v}{1.24}\bigr), & v<0.62,\\[2pt]
    W\sqrt{1-q^2},\quad q=\frac{v-0.62}{0.38}, & v\ge 0.62,
  \end{cases}
  \qquad W\in[0.40,0.48],
\]
a narrow claw, widest at $0.62$, with an elliptical cap. The margin is moved along the
normal by $0.03\,v\,n(6v)$, where $n$ is a three-octave gradient noise seeded per petal, so
it wavers more toward the tip; four petals in ten also get a shallow notch. The outline is
sampled once and used twice: as a clipping path, and as the stroke that gives the rim.
Inside the clip the petal is painted like a watercolor, in the order a painter would: a
flat wash at $\operatorname{tone}(t_0)$, $t_0\in[0.30,0.46]$; eight stacked white blobs at
the base, each at $10\%$ opacity, for the pale heart; fifteen irregular blotches
multiplied over the wash, most darker, some lighter, and a dozen dark specks for the
granulation; five to eight darker brush strokes running out from the base; the outline
stroked $1.6\,\ell\,$bp wide for a petal $\ell\,$cm long and multiplied, of which only the inner half shows, for the pigment
that pools at the edge of a wash; nine to thirteen white streaks lifted out in a fan from
the base; and a few fine dark veins. There is no outline in the usual sense.
\end{construction}

\begin{construction}[Tapered strokes]
Every streak, vein and rib is a filled shape rather than a stroke, so that it tapers. For a
cubic B\'ezier with control points $P_0,\dots,P_3$ and widths $w_0$, $w_1$ at its ends, push
each $P_i$ along its normal by $\pm\frac12 w(i/3)$, with the normals at $P_1$ and $P_2$ taken
from the chords $P_2-P_0$ and $P_3-P_1$. The two offset curves joined at the ends are a
tapered stroke in two B\'ezier segments, five times smaller in the file than a sampled
ribbon and indistinguishable from one for a gentle curve.
\end{construction}

\begin{construction}[Flowers]
A \emph{flower} is five petals at $72^\circ$ intervals with $\pm6^\circ$ of jitter and lengths
in $[0.9,1.06]$, then a heart of ten white rays, a white disc, fourteen stamens and a copper
eye. A quarter of the flowers are painted a shade deeper. A \emph{lily} is six pointed
petals, $w(v)=2W\sin^{0.85}(\pi v^{0.8})(1-0.3v)$, arranged in a star and then squashed
vertically by a factor in $[0.45,0.62]$: the star is seen at a tilt, so the petals that point
away are short, and since they are drawn first and darker they read as the back of the
flower. A \emph{blossom} is the same idea at a few millimeters, and a \emph{floret}, the
smallest, is five petals filled as one shape around a dark eye; below $0.7\,$mm in radius it is a
ten-pointed star. Buds, bellflowers and sepals are blades on the same spine.
\end{construction}

\begin{construction}[Leaves]
A leaf of length $1$ has the envelope $e(v)=2.3\,W v^{0.55}(1-v)^{0.95}$. An entire leaf is
that envelope with fine teeth. A lobed leaf, like acanthus, is a narrow body of half-width
$(1-0.75\,d)\,e(v)$, where $d$ is the depth of the cut, together with three to five pointed
lobes on each side; each lobe is itself a blade that leaves the midrib at
$38^\circ$--$58^\circ$ toward the tip and reaches the envelope. Body and lobes are filled as
one shape, under the nonzero rule, so there is no seam. One half of the leaf is multiplied
darker, as if it were folded along its rib; then come the blotches, the pooled edge, and a
white midrib with a white vein along every lobe.
\end{construction}

\begin{construction}[Sprays]
A \emph{cyme}, the inflorescence of the forget-me-not, is a stem whose heading turns by
$\pm\tau(s/L)^{2.3}$ at arc length $s$, with $\tau\in[170^\circ,270^\circ]$: straight at first,
it curls like a scorpion's tail at the end. Eleven to seventeen florets crowd along it,
shrinking into buds toward the curl, and a second, shorter curl forks off it three times in
five. A \emph{sprig} is a stem that forks two or three ways four times, each fork $0.58$ to
$0.74$ the length of its parent, with florets and buds at the tips. An \emph{umbel} fans nine
to thirteen rays over $110^\circ$, each ending in a small umbel of starry florets.
\end{construction}

\subsubsection{Composition}

The garden is a plant, not an arrangement. It grows from one root in the bottom-left
corner of the page, where the base of the stems curls into a scroll. Three stems run up the
column through hand-picked control points, interpolated by Catmull--Rom splines and jittered
by a millimeter or so on each page: the tallest ends in the large flower at the top, one in
a lily, one in a flower half-way up. A fourth stem leaves the root along the bottom of the
page, carrying small leaves that thin out toward its end, and runs out into a copper
hairline; it is the footer. The header is a printer's ivy leaf in copper whose stem becomes
the head rule.

Leaves hang at chosen fractions of the stems' length, alternately left and right. Every
other motif sits at a chosen point and is joined to a chosen stem by a \emph{branch}: a cubic
B\'ezier that leaves the stem along the stem's own direction turned $28^\circ$ to $48^\circ$
toward the target, and arrives along the motif's axis. Because a branch starts tangent to
its parent and ends tangent to its flower, every flower can be traced back to the root, and
that, more than any detail of the painting, is what makes the margin look grown. Behind
everything are one or two \emph{ghosts}: flowers painted with every tone lifted toward the
paper, $t\mapsto t+0.62(1-t)$, and without their fine detail, which give the column depth.

\begin{remark}[Procedural is not random]
The generator is seeded from the page number before the garden is drawn. Page $n$
therefore receives the same garden on every compilation, and a different one from page
$n+1$. The variation the reader notices between pages is real; the reproducibility is what
makes the document stable enough to typeset.
\end{remark}

\begin{proposition}
Let $P$ be the number of primitive instances on a page and $S$ the mean number of path
points per instance. The cost of a page is $\Theta(PS)$ floating-point operations in Lua and
$\Theta(PS)$ bytes of output, and \TeX\ does no arithmetic at all; in practice a page takes
well under a second, for $P\approx 400$ and $S\approx 100$.
\end{proposition}

\begin{proof}[Proof sketch]
Every coordinate is computed once in Lua, transformed to the page, and printed with one
decimal of a PDF point into a single literal node, which \TeX\ only carries to the page.
The same curves can be computed in \TeX\ itself with \texttt{pgfmath}, but it does
arithmetic by macro expansion in \TeX\ dimensions, and a page then takes about a minute. The
size of the output is the real constraint: about
three megabytes of path operators per page, which compress to under one.
\end{proof}

\begin{exercise}
The composition above is authored. Replace it by a growth process in the spirit of Wong,
Zongker and Salesin (\emph{Computer-generated floral ornament}, SIGGRAPH 1998): starting
from the stems, repeatedly find the largest free disc in the margin that is within reach of
a stem, choose a primitive that fits it, and grow a branch to it, until no disc larger than
a threshold remains. The ornament then adapts to any text block, and a page whose text
ends early grows a garden into the space below it.
\end{exercise}
